%! Parametrization of Implicitly Defined Functions — Unit 4, Lesson 4.7 — Lesson Plan
\documentclass[10pt]{article}
\usepackage{precalculus-article}
\usepackage{precalculus-boxes}
\usepackage{graphicx}
\graphicspath{{images/}}
\newcommand{\TallMath}[1]{$\displaystyle #1\rule[-1.4em]{0pt}{3.2em}$}
\newcommand{\CourseName}{Precalculus}
\newcommand{\SchoolYear}{2026--2027}
\newcommand{\MeetingLength}{55 minutes}
\newcommand{\UnitNumberName}{Unit 4: Functions Involving Parameters, Vectors, and Matrices \quad}
\newcommand{\LessonNumberName}{Lesson 4.7: Parametrization of Implicitly Defined Functions}

\begin{document}

%----------------------------------------------------------------------
% Title Block
%----------------------------------------------------------------------
\begin{center}
  {\Large\bfseries \CourseName: \SchoolYear} \\[4pt]
  {\small\bfseries \UnitNumberName \LessonNumberName \quad (2 instructional periods, \MeetingLength\ each)}
\end{center}

%----------------------------------------------------------------------
% Primary Objective
%----------------------------------------------------------------------
\begin{tcolorbox}[colback=sky, colframe=navy, boxrule=0.9pt, arc=2mm,
    left=2mm, right=2mm, top=1.5mm, bottom=1.5mm]
  \textbf{Primary Objective:} Students will express implicitly defined
  curves --- ellipses, circles, and parabolas --- as parametric functions
  using trigonometric parametrizations and algebraic substitutions, and
  apply parametric analysis (extrema from component functions, axis
  intercepts from zeros) to these curves. \textit{(Big Idea: Functions)}
\end{tcolorbox}

\vspace{0.08in}

%----------------------------------------------------------------------
% Priority Ideas & Skills
%----------------------------------------------------------------------
\begin{skillbox}[Priority Ideas \& Skills]{goldbox}
\begin{tabularx}{\linewidth}{@{} p{0.46\linewidth} X @{}}
\textbf{AP Skills} & \textbf{Key Understandings (EKs)} \\
\midrule
\textbf{1.B} \textit{Express in equivalent forms} ---
  Rewrite an implicitly defined equation (circle, ellipse, parabola)
  as a pair of parametric component functions that trace the curve,
  choosing a parameter substitution that makes the implicit equation
  a Pythagorean identity or a direct algebraic substitution.
&
\textbf{EK 4.7.A.1}: An implicitly defined equation in two variables
can often be expressed as a parametric function by introducing a
parameter $t$ that traces the curve. \\[4pt]

\textbf{2.A} \textit{Identify information from representations} ---
  Given a parametric pair $(x(t), y(t))$ derived from an implicit
  curve, extract the extrema of $x$ and $y$, the axis intercepts, and
  the direction of motion as $t$ increases.
&
\textbf{EK 4.7.A.2}: For an ellipse
$\tfrac{(x-h)^2}{a^2}+\tfrac{(y-k)^2}{b^2}=1$, the parametrization
$x(t)=h+a\cos t$, $y(t)=k+b\sin t$, $0\le t\le 2\pi$, traces the
ellipse counterclockwise. \\[4pt]

&
\textbf{EK 4.7.A.3}: For a parabola $y-k=a(x-h)^2$, one
parametrization is $x(t)=h+t$, $y(t)=k+at^2$ for $t\in\mathbb{R}$
(or a restricted domain). \\[4pt]

&
\textbf{EK 4.7.A.4}: Once parametrized, parametric analysis tools
(direction of motion, extrema, intercepts) can be applied to curves
that are not functions in the Cartesian sense. \\
\end{tabularx}
\end{skillbox}

\vspace{0.08in}

%----------------------------------------------------------------------
% Vocabulary
%----------------------------------------------------------------------
\begin{skillbox}[{Vocabulary, Concepts, \& Theorems}]{greenbox}
\renewcommand{\arraystretch}{1.4}
\begin{tabularx}{\linewidth}{@{} p{0.26\linewidth} X @{}}
\textbf{Term} & \textbf{Definition / Note} \\
\midrule
implicit curve &
  A curve defined by an equation relating $x$ and $y$ that cannot
  always be written as a single explicit function $y=f(x)$;
  e.g.\ $\tfrac{x^2}{9}+\tfrac{y^2}{4}=1$ (ellipse). \\[4pt]
parametrization &
  A pair of component functions $(x(t), y(t))$ that traces an implicit
  curve as $t$ varies; the parameter $t$ may represent time or another
  quantity. \\[4pt]
ellipse parametrization &
  \TallMath{x(t)=h+a\cos t,\; y(t)=k+b\sin t,\; 0\le t\le 2\pi;}
  derived from the Pythagorean identity $\cos^2 t+\sin^2 t=1$. \\[4pt]
parabola parametrization &
  For $y-k=a(x-h)^2$: set $x(t)=h+t$, $y(t)=k+at^2$, $t\in\mathbb{R}$;
  the vertex occurs at $t=0$, giving $(h, k)$. \\
\end{tabularx}
\end{skillbox}

\vspace{0.08in}

%----------------------------------------------------------------------
% Activate Prior Knowledge
%----------------------------------------------------------------------
\begin{skillbox}[Activate Prior Knowledge \& Spiral Review]{sky}
\begin{tabularx}{\linewidth}{@{}X c@{}}
\begin{minipage}[t]{0.96\linewidth}\vspace{0pt}
\textbf{Spiral topics activated during Warm-Up (first 8 minutes):}
\begin{itemize}[leftmargin=1.4em, itemsep=2pt]
  \item \textbf{Unit circle parametrization} (Lesson 4.4):
    $(x(t),y(t))=(\cos t,\sin t)$ for $0\le t\le 2\pi$ --- the foundation
    for all trig-based parametrizations.
  \item \textbf{Ellipse intercepts from standard form} (Lesson 4.6):
    reading $a$ and $b$ from $\tfrac{x^2}{a^2}+\tfrac{y^2}{b^2}=1$ to
    locate $x$- and $y$-intercepts.
  \item \textbf{Evaluating parametric component functions at specific $t$}
    (Lesson 4.2): computing $(x(t), y(t))$ at given values and checking
    whether those points satisfy an implicit equation.
\end{itemize}
\end{minipage}
&
\end{tabularx}
\end{skillbox}

\vspace{0.08in}

%----------------------------------------------------------------------
% Hook
%----------------------------------------------------------------------
\begin{skillbox}[Hook --- Entry Event]{sky}
\textbf{Scenario:} The ellipse $\dfrac{x^2}{9}+\dfrac{y^2}{4}=1$ is
\emph{not} a function --- it fails the vertical line test. But what if
a particle traces this ellipse over time? Then we can ask:
\textit{Where is the particle at $t=\pi/3$? When does it cross the
$y$-axis? What is its leftmost position?} These are parametric questions,
and answering them requires converting the implicit equation into a
parametric form.

\textbf{Discussion prompts:}
\begin{enumerate}[leftmargin=1.4em, itemsep=3pt]
  \item \textit{``The identity $\cos^2 t + \sin^2 t = 1$ looks a lot like
    $\tfrac{x^2}{9}+\tfrac{y^2}{4}=1$ if you write $x/3 = \cos t$ and
    $y/2 = \sin t$. How could you use that observation?''}
  \item \textit{``Once you have $(x(t),y(t))$, how do you find when the
    particle reaches its rightmost point?''}
\end{enumerate}

\textbf{Key tension:} Implicit equations describe the \emph{shape};
parametric forms add \emph{motion} --- direction, position at time $t$,
and extrema using tools already developed in Lessons 4.2 and 4.3.
\end{skillbox}

\vspace{0.08in}

%----------------------------------------------------------------------
% Lesson
%----------------------------------------------------------------------
\begin{skillbox}[Lesson --- Instruction Sequence (Period 1)]{sky}
\begin{multicols}{2}

\textbf{Step 1 --- Review: Unit Circle and Scaling} (6 min)

Revisit $(x(t),y(t))=(\cos t, \sin t)$, $0\le t\le 2\pi$.
Key identity: $\cos^2 t+\sin^2 t=1$.
Scaling: if $x=a\cos t$ and $y=b\sin t$, substitute:
$\tfrac{x^2}{a^2}+\tfrac{y^2}{b^2}=\cos^2 t+\sin^2 t=1$.
Ask: \textit{``When $a\ne b$, what shape do we trace?''}
An ellipse. This is the core idea of EK 4.7.A.1.

\columnbreak

\textbf{Step 2 --- Ellipse Parametrization} (15 min)

Standard form: $\tfrac{(x-h)^2}{a^2}+\tfrac{(y-k)^2}{b^2}=1$.
Let $\tfrac{x-h}{a}=\cos t$ and $\tfrac{y-k}{b}=\sin t$, so:
\[
  x(t)=h+a\cos t,\quad y(t)=k+b\sin t,\quad 0\le t\le 2\pi.
\]
Verify: $\tfrac{(h+a\cos t - h)^2}{a^2}+\tfrac{(k+b\sin t - k)^2}{b^2}
= \cos^2 t+\sin^2 t=1.\;\checkmark$
At $t=0$: $(h+a,\;k)$ --- rightmost. At $t=\pi/2$: $(h,\;k+b)$ --- top.

Work through: $\tfrac{x^2}{16}+\tfrac{y^2}{4}=1$.
$h=0,\,k=0,\,a=4,\,b=2$.
Parametrization: $x(t)=4\cos t$, $y(t)=2\sin t$, $0\le t\le 2\pi$.
(EK 4.7.A.2)

\end{multicols}

\begin{multicols}{2}

\textbf{Step 3 --- Parabola Parametrization} (12 min)

For $y-k=a(x-h)^2$: $x$ is unrestricted, so set $x(t)=h+t$, giving
$y(t)=k+at^2$, $t\in\mathbb{R}$. Vertex at $t=0$: $(h,k)$.
For a sideways parabola $x-h=a(y-k)^2$: set $y(t)=k+t$, $x(t)=h+at^2$.

Example: $y=x^2-4x+3=(x-2)^2-1$. Vertex $(2,-1)$, $a=1$.
Parametrization: $x(t)=2+t$, $y(t)=-1+t^2$.
Evaluate at $t=-2,-1,0,1,2$: $x = 0,1,2,3,4$ and $y=3,0,-1,0,3$.
(EK 4.7.A.3)

\columnbreak

\textbf{Step 4 --- Apply Parametric Analysis} (10 min)

For the ellipse $x(t)=3\cos t$, $y(t)=5\sin t$:
\begin{itemize}[leftmargin=1.2em, itemsep=2pt]
  \item Rightmost point: $x(t)=3$ when $t=0$ $\Rightarrow$ $(3,0)$.
  \item Leftmost point: $x(t)=-3$ when $t=\pi$ $\Rightarrow$ $(-3,0)$.
  \item $y$-axis crossing ($x(t)=0$): $\cos t=0$, so $t=\pi/2,\;3\pi/2$.
    Points: $(0,5)$ and $(0,-5)$.
  \item $x$-axis crossing ($y(t)=0$): $\sin t=0$, so $t=0,\;\pi$.
    Points: $(3,0)$ and $(-3,0)$.
\end{itemize}
Reinforce: EK 4.7.A.4 --- parametric tools apply to non-function curves.

\end{multicols}

\smallskip
\textbf{Pacing note:} Steps 1--2 fill Period 1; Steps 3--4 open Period 2 before activity.
\end{skillbox}

\vspace{0.08in}

%----------------------------------------------------------------------
% Active Monitoring
%----------------------------------------------------------------------
\begin{skillbox}[Active Monitoring]{redbox}
\begin{itemize}[leftmargin=1.4em, itemsep=3pt]
  \item \textbf{Mixing up $a$ and $b$}: Students may put $b$ under $x$
    and $a$ under $y$. Prompt: \textit{``Which denominator appears below
    $(x-h)^2$ in standard form? That value is $a^2$, so $a$ goes with $x$.''}
  \item \textbf{Forgetting the center shift}: Students write $x(t)=a\cos t$
    even when $h\ne 0$. Ask: \textit{``At $t=0$, your formula gives $x=a$.
    What should the rightmost point of this ellipse be?''}
  \item \textbf{Domain for parabola}: Students may restrict $t$ to $[0,2\pi]$
    by analogy with ellipses. Clarify: for parabolas, $t$ runs over all
    reals (the curve extends without bound).
  \item \textbf{Extrema by picture vs.\ component function}: Students read
    extrema from the Cartesian graph rather than solving the component
    equation. Reinforce: set $x(t)$ equal to its maximum value and solve
    for $t$; same for $y(t)$.
\end{itemize}
\end{skillbox}

\vspace{0.08in}

%----------------------------------------------------------------------
% Group Work & Differentiation
%----------------------------------------------------------------------
\begin{skillbox}[Group Work \& Differentiation]{redbox}
Students work in groups of 3--4 on the Mission Control activity (tiered).

\begin{itemize}[leftmargin=1.4em, itemsep=4pt]
  \item \textbf{Tier R (Remediate)}: Given the ellipse
    $\tfrac{x^2}{16}+\tfrac{y^2}{9}=1$, verify that $x(t)=4\cos t$,
    $y(t)=3\sin t$ satisfies the equation at $t=0$, $t=\pi/2$, and $t=\pi$.
  \item \textbf{Tier A (Approaching Proficiency)}: Write a parametrization
    for $\tfrac{(x-2)^2}{25}+\tfrac{(y+1)^2}{9}=1$ and find the
    rightmost and leftmost extreme positions using the parametrization.
  \item \textbf{Tier E (Extension)}: For the sideways parabola
    $x-1=2(y+3)^2$, write a parametrization with $y(t)=-3+t$, identify
    the vertex coordinates, and find the axis intercepts by solving
    $x(t)=0$ and $y(t)=0$.
\end{itemize}
\end{skillbox}

\vspace{0.08in}

%----------------------------------------------------------------------
% Individual Work & Assessment
%----------------------------------------------------------------------
\begin{skillbox}[Individual Work \& Assessment]{redbox}
Exit ticket administered independently (final 5 minutes of Period 2).

\begin{enumerate}[leftmargin=1.4em, itemsep=4pt]
  \item Write a parametrization for
    $\dfrac{(x+1)^2}{9}+\dfrac{(y-2)^2}{16}=1$.
  \item Using your parametrization, find the coordinates of the topmost
    and bottommost points on the ellipse.
\end{enumerate}

\textbf{Data use:} Students who cannot perform the trig substitution
without prompting need a reference card linking $\cos^2 t+\sin^2 t=1$
to ellipse standard form; revisit at the start of the next lesson.
Students who find extrema by returning to the Cartesian picture rather
than solving the component equation need additional parametric analysis
practice (Lesson 4.2 review).
\end{skillbox}

\vspace{0.08in}

%----------------------------------------------------------------------
% Reinforcement & Extension
%----------------------------------------------------------------------
\begin{skillbox}[Reinforcement \& Extension]{goldbox}
\textbf{Homework (Lesson 4.7):} 5--6 problems covering: parametrizing
ellipses and circles in standard form, parametrizing parabolas by
substitution, finding extrema and intercepts from the parametric form,
and one AP-style multiple-choice item comparing two valid parametrizations
of the same ellipse.

\textbf{Extension:} For the ellipse $x(t)=a\cos t$, $y(t)=b\sin t$,
what substitution traces the ellipse \emph{clockwise}? How does
replacing $t$ with $-t$ affect $\cos t$ and $\sin t$ separately?

\textbf{Preview of Lesson 4.8:} Ask students:
\textit{``From $x=h+t$, $y=k+at^2$, can you eliminate $t$ to recover
a Cartesian equation? What equation in $x$ and $y$ do you get?''} This
reversal --- parametric back to implicit --- will be the focus next.
\end{skillbox}

\end{document}
